\chapter{Mixed Structural Table Example}

\section{Purpose}

This chapter documents a compact structural example that exercises the main
architectural decisions developed in the preceding chapters:
\begin{itemize}
  \item a heterogeneous structural model stored through
        \texttt{StructuralElement},
  \item formulation-specific reconstruction through
        \texttt{StructuralReductionPolicy<ElementT>},
  \item multi-block VTK export through \texttt{StructuralVTMExporter},
  \item coexistence of beam and shell elements on a shared geometric support.
\end{itemize}

The example is intentionally simple enough to inspect by hand, but rich enough
to validate the polymorphic structural pipeline end to end.

\section{Physical Model}

The structure resembles a table:
\begin{itemize}
  \item four corner columns,
  \item four perimeter beams on the top frame,
  \item one shell slab spanning the frame.
\end{itemize}

\subsection{Geometry}

The adopted dimensions are:
\[
L = 5.0, \qquad H = 3.0, \qquad t_\mathrm{slab} = 0.18.
\]

The four column bases are placed at:
\[
(0,0,0),\; (L,0,0),\; (L,L,0),\; (0,L,0),
\]
and the top frame lies in the plane:
\[
z = H.
\]

The slab is discretised as a structured \(4\times 4\) MITC4 mesh over the
square domain \([0,L]\times[0,L]\).

\subsection{Sections and Materials}

The example uses two distinct beam formulations and one shell formulation:
\begin{itemize}
  \item \textbf{Columns}: \texttt{BeamElement<TimoshenkoBeam3D,3>} with
        rectangular section \(0.40\times 0.40\), Young's modulus
        \(E_c = 21000\), Poisson ratio \(\nu = 0.20\).
  \item \textbf{Perimeter beams}: \texttt{TimoshenkoBeamN<2>} with
        rectangular section \(0.60\times 0.40\), Young's modulus
        \(E_b = 24000\), Poisson ratio \(\nu = 0.20\).
  \item \textbf{Slab}: \texttt{ShellElement<MindlinReissnerShell3D>} with
        thickness \(t = 0.18\), Young's modulus \(E_s = 26000\), Poisson ratio
        \(\nu = 0.20\).
\end{itemize}

The shear modulus is reconstructed from isotropic elasticity:
\[
G = \frac{E}{2(1+\nu)}.
\]

For rectangular beams the section properties are computed directly from the
section dimensions:
\[
A = b h,
\qquad
I_y = \frac{b h^3}{12},
\qquad
I_z = \frac{h b^3}{12},
\]
and the Saint-Venant torsional constant is approximated by
\[
J \approx \frac{b_\mathrm{min}^3 h_\mathrm{max}}{3}
\left(1 - 0.63 \frac{b_\mathrm{min}}{h_\mathrm{max}}\right).
\]

\section{Loads and Boundary Conditions}

All nodes with \(z=0\) are fully clamped through the structural model
interface.

Two load cases are superposed:
\begin{enumerate}
  \item a uniform slab pressure
        \[
        p_z = -0.12,
        \]
        converted to consistent nodal forces by direct integration of the shell
        interpolation,
        \[
        \mathbf{f}_I = \int_{\Omega_e} N_I \mathbf{t}\, dA,
        \qquad
        \mathbf{t} = (0,0,p_z)^T;
        \]
  \item a lateral point force applied at the corner node \((L,L,H)\),
        \[
        F_x = 0.20.
        \]
\end{enumerate}

This combination is deliberately chosen to excite:
\begin{itemize}
  \item vertical flexure of the slab,
  \item frame action in the perimeter beams,
  \item bending plus sway of the full table.
\end{itemize}

\section{Architectural Realisation}

\subsection{Why the Domain Contains Only the Slab}

The current \texttt{Domain} layer is intentionally aligned with a pure
\texttt{DMPlex} surface/volume topology.  The slab therefore lives inside the
PETSc mesh, whereas the columns and beams are introduced as external line
geometries bound to the same \texttt{Node}/\texttt{Vertex} storage.  This
keeps the domain ontologically clean:
\begin{itemize}
  \item the 2D topological mesh remains a shell/slab mesh,
  \item the line members remain structural geometries rather than fake 2D
        entities,
  \item the mixed structural model is assembled one layer above the domain.
\end{itemize}

\subsection{Polymorphic Structural Container}

The final element container is
\[
\texttt{std::vector<StructuralElement>},
\]
owned by
\[
\texttt{Model<TimoshenkoBeam3D,\; SmallStrain,\; 6,\; SingleElementPolicy<StructuralElement>>}.
\]

This means that the analysis sees one structural family with a common
finite-element interface, while post-processing can still recover the concrete
type through the structural type-erasure layer:
\begin{itemize}
  \item \texttt{BeamElement<TimoshenkoBeam3D,3>} for columns,
  \item \texttt{TimoshenkoBeamN<2>} for edge beams,
  \item \texttt{ShellElement<MindlinReissnerShell3D>} for slab cells.
\end{itemize}

The example therefore validates the intended design point: runtime
heterogeneity at the model boundary, formulation-specific compile-time
reconstruction inside each policy.

\section{Quadrature Contract for \texttt{TimoshenkoBeamN}}

The example uncovered an architectural inconsistency in
\texttt{TimoshenkoBeamN<N>}.  The element uses a reduced constitutive
integration with
\[
n_\mathrm{gp} = N - 1,
\]
but its public topology query originally forwarded
\texttt{geometry\_->num\_integration\_points()}, which depends on the
underlying geometry integrator.  For \texttt{TimoshenkoBeamN<2>} this created
an inconsistent state:
\begin{itemize}
  \item the element stored one constitutive section,
  \item the exporter believed that the element had two Gauss points when the
        geometry used \texttt{GaussLegendreCellIntegrator<2>}.
\end{itemize}

The fix adopted here is structural, not cosmetic:
\begin{enumerate}
  \item \texttt{TimoshenkoBeamN<N>::num\_integration\_points()} now returns
        the reduced value \(N-1\).
  \item The constructor now enforces the quadrature contract:
        the attached geometry must also expose exactly \(N-1\) integration
        points.
  \item The table example therefore assigns
        \texttt{GaussLegendreCellIntegrator<1>} to the perimeter beams and
        \texttt{GaussLegendreCellIntegrator<2>} to the columns.
\end{enumerate}

This is important beyond visualisation.  The reconstruction layer and any
future scale-transfer operator must reason about the actual constitutive sites
of the formulation, not about an unrelated geometric quadrature.

\section{VTK Output}

The structural exporter writes a \texttt{.vtm} index together with one
\texttt{.vtp} sidecar per semantic block:
\begin{itemize}
  \item \texttt{table\_frame\_structural.vtm},
  \item \texttt{table\_frame\_structural\_axis.vtp},
  \item \texttt{table\_frame\_structural\_surface.vtp},
  \item \texttt{table\_frame\_structural\_sections.vtp},
  \item \texttt{table\_frame\_structural\_material\_sites.vtp},
  \item \texttt{table\_frame\_structural\_fibers.vtp}.
\end{itemize}

These blocks have a fixed semantic meaning:
\begin{description}
  \item[\texttt{axis}] beam centre-lines and shell boundary carriers;
  \item[\texttt{surface}] reconstructed beam skins and shell top/bottom faces;
  \item[\texttt{sections}] section stations and shell through-thickness lines;
  \item[\texttt{material\_sites}] constitutive sampling sites;
  \item[\texttt{fibers}] fiber carriers when a fiber section is present.
\end{description}

The split into \texttt{.vtp} sidecars is intentional.  It avoids the fragile
writer path observed with the monolithic multi-block writer while preserving a
semantically clean \texttt{.vtm} interface for ParaView.

\section{Observed Structural Response}

For the current parameter set, the example run stored in the repository
reported:
\[
\max |u| = 1.3651\times 10^{-2}, \qquad
\max |u_x| = 5.30\times 10^{-4}, \qquad
\max |u_z| = 1.3651\times 10^{-2}.
\]

These values are not presented as a benchmark target.  Their role here is to
confirm that:
\begin{itemize}
  \item the mixed structural solve completes,
  \item the reconstructed carriers can be exported without crashing,
  \item the example exercises both vertical slab deformation and lateral frame
        sway.
\end{itemize}

\section{Code Path}

The complete example is implemented in \texttt{main.cpp} as
\texttt{run\_structural\_table\_example()}.  The relevant pieces are:
\begin{itemize}
  \item geometric construction of slab, columns, and beams;
  \item consistent shell pressure assembly through
        \texttt{apply\_uniform\_shell\_surface\_load()};
  \item polymorphic structural model assembly through
        \texttt{StructuralElement};
  \item multi-block structural export through
        \texttt{export\_structural\_vtm()}.
\end{itemize}

This example now serves as the reference mixed structural demonstration for the
current architecture.
